% Options for packages loaded elsewhere
% Options for packages loaded elsewhere
\PassOptionsToPackage{unicode}{hyperref}
\PassOptionsToPackage{hyphens}{url}
\PassOptionsToPackage{dvipsnames,svgnames,x11names}{xcolor}
%
\documentclass[
  10pt,
  letterpaper,
  DIV=11,
  numbers=noendperiod]{scrartcl}
\usepackage{xcolor}
\usepackage[top=20mm,bottom=20mm,left=20mm,right=20mm]{geometry}
\usepackage{amsmath,amssymb}
\setcounter{secnumdepth}{5}
\usepackage{iftex}
\ifPDFTeX
  \usepackage[T1]{fontenc}
  \usepackage[utf8]{inputenc}
  \usepackage{textcomp} % provide euro and other symbols
\else % if luatex or xetex
  \usepackage{unicode-math} % this also loads fontspec
  \defaultfontfeatures{Scale=MatchLowercase}
  \defaultfontfeatures[\rmfamily]{Ligatures=TeX,Scale=1}
\fi
\usepackage{lmodern}
\ifPDFTeX\else
  % xetex/luatex font selection
\fi
% Use upquote if available, for straight quotes in verbatim environments
\IfFileExists{upquote.sty}{\usepackage{upquote}}{}
\IfFileExists{microtype.sty}{% use microtype if available
  \usepackage[]{microtype}
  \UseMicrotypeSet[protrusion]{basicmath} % disable protrusion for tt fonts
}{}
\makeatletter
\@ifundefined{KOMAClassName}{% if non-KOMA class
  \IfFileExists{parskip.sty}{%
    \usepackage{parskip}
  }{% else
    \setlength{\parindent}{0pt}
    \setlength{\parskip}{6pt plus 2pt minus 1pt}}
}{% if KOMA class
  \KOMAoptions{parskip=half}}
\makeatother
% Make \paragraph and \subparagraph free-standing
\makeatletter
\ifx\paragraph\undefined\else
  \let\oldparagraph\paragraph
  \renewcommand{\paragraph}{
    \@ifstar
      \xxxParagraphStar
      \xxxParagraphNoStar
  }
  \newcommand{\xxxParagraphStar}[1]{\oldparagraph*{#1}\mbox{}}
  \newcommand{\xxxParagraphNoStar}[1]{\oldparagraph{#1}\mbox{}}
\fi
\ifx\subparagraph\undefined\else
  \let\oldsubparagraph\subparagraph
  \renewcommand{\subparagraph}{
    \@ifstar
      \xxxSubParagraphStar
      \xxxSubParagraphNoStar
  }
  \newcommand{\xxxSubParagraphStar}[1]{\oldsubparagraph*{#1}\mbox{}}
  \newcommand{\xxxSubParagraphNoStar}[1]{\oldsubparagraph{#1}\mbox{}}
\fi
\makeatother

\usepackage{color}
\usepackage{fancyvrb}
\newcommand{\VerbBar}{|}
\newcommand{\VERB}{\Verb[commandchars=\\\{\}]}
\DefineVerbatimEnvironment{Highlighting}{Verbatim}{commandchars=\\\{\}}
% Add ',fontsize=\small' for more characters per line
\usepackage{framed}
\definecolor{shadecolor}{RGB}{241,243,245}
\newenvironment{Shaded}{\begin{snugshade}}{\end{snugshade}}
\newcommand{\AlertTok}[1]{\textcolor[rgb]{0.68,0.00,0.00}{#1}}
\newcommand{\AnnotationTok}[1]{\textcolor[rgb]{0.37,0.37,0.37}{#1}}
\newcommand{\AttributeTok}[1]{\textcolor[rgb]{0.40,0.45,0.13}{#1}}
\newcommand{\BaseNTok}[1]{\textcolor[rgb]{0.68,0.00,0.00}{#1}}
\newcommand{\BuiltInTok}[1]{\textcolor[rgb]{0.00,0.23,0.31}{#1}}
\newcommand{\CharTok}[1]{\textcolor[rgb]{0.13,0.47,0.30}{#1}}
\newcommand{\CommentTok}[1]{\textcolor[rgb]{0.37,0.37,0.37}{#1}}
\newcommand{\CommentVarTok}[1]{\textcolor[rgb]{0.37,0.37,0.37}{\textit{#1}}}
\newcommand{\ConstantTok}[1]{\textcolor[rgb]{0.56,0.35,0.01}{#1}}
\newcommand{\ControlFlowTok}[1]{\textcolor[rgb]{0.00,0.23,0.31}{\textbf{#1}}}
\newcommand{\DataTypeTok}[1]{\textcolor[rgb]{0.68,0.00,0.00}{#1}}
\newcommand{\DecValTok}[1]{\textcolor[rgb]{0.68,0.00,0.00}{#1}}
\newcommand{\DocumentationTok}[1]{\textcolor[rgb]{0.37,0.37,0.37}{\textit{#1}}}
\newcommand{\ErrorTok}[1]{\textcolor[rgb]{0.68,0.00,0.00}{#1}}
\newcommand{\ExtensionTok}[1]{\textcolor[rgb]{0.00,0.23,0.31}{#1}}
\newcommand{\FloatTok}[1]{\textcolor[rgb]{0.68,0.00,0.00}{#1}}
\newcommand{\FunctionTok}[1]{\textcolor[rgb]{0.28,0.35,0.67}{#1}}
\newcommand{\ImportTok}[1]{\textcolor[rgb]{0.00,0.46,0.62}{#1}}
\newcommand{\InformationTok}[1]{\textcolor[rgb]{0.37,0.37,0.37}{#1}}
\newcommand{\KeywordTok}[1]{\textcolor[rgb]{0.00,0.23,0.31}{\textbf{#1}}}
\newcommand{\NormalTok}[1]{\textcolor[rgb]{0.00,0.23,0.31}{#1}}
\newcommand{\OperatorTok}[1]{\textcolor[rgb]{0.37,0.37,0.37}{#1}}
\newcommand{\OtherTok}[1]{\textcolor[rgb]{0.00,0.23,0.31}{#1}}
\newcommand{\PreprocessorTok}[1]{\textcolor[rgb]{0.68,0.00,0.00}{#1}}
\newcommand{\RegionMarkerTok}[1]{\textcolor[rgb]{0.00,0.23,0.31}{#1}}
\newcommand{\SpecialCharTok}[1]{\textcolor[rgb]{0.37,0.37,0.37}{#1}}
\newcommand{\SpecialStringTok}[1]{\textcolor[rgb]{0.13,0.47,0.30}{#1}}
\newcommand{\StringTok}[1]{\textcolor[rgb]{0.13,0.47,0.30}{#1}}
\newcommand{\VariableTok}[1]{\textcolor[rgb]{0.07,0.07,0.07}{#1}}
\newcommand{\VerbatimStringTok}[1]{\textcolor[rgb]{0.13,0.47,0.30}{#1}}
\newcommand{\WarningTok}[1]{\textcolor[rgb]{0.37,0.37,0.37}{\textit{#1}}}

\usepackage{longtable,booktabs,array}
\usepackage{calc} % for calculating minipage widths
% Correct order of tables after \paragraph or \subparagraph
\usepackage{etoolbox}
\makeatletter
\patchcmd\longtable{\par}{\if@noskipsec\mbox{}\fi\par}{}{}
\makeatother
% Allow footnotes in longtable head/foot
\IfFileExists{footnotehyper.sty}{\usepackage{footnotehyper}}{\usepackage{footnote}}
\makesavenoteenv{longtable}
\usepackage{graphicx}
\makeatletter
\newsavebox\pandoc@box
\newcommand*\pandocbounded[1]{% scales image to fit in text height/width
  \sbox\pandoc@box{#1}%
  \Gscale@div\@tempa{\textheight}{\dimexpr\ht\pandoc@box+\dp\pandoc@box\relax}%
  \Gscale@div\@tempb{\linewidth}{\wd\pandoc@box}%
  \ifdim\@tempb\p@<\@tempa\p@\let\@tempa\@tempb\fi% select the smaller of both
  \ifdim\@tempa\p@<\p@\scalebox{\@tempa}{\usebox\pandoc@box}%
  \else\usebox{\pandoc@box}%
  \fi%
}
% Set default figure placement to htbp
\def\fps@figure{htbp}
\makeatother





\setlength{\emergencystretch}{3em} % prevent overfull lines

\providecommand{\tightlist}{%
  \setlength{\itemsep}{0pt}\setlength{\parskip}{0pt}}



 


\usepackage{booktabs}
\usepackage{longtable}
\usepackage{array}
\usepackage{multirow}
\usepackage{wrapfig}
\usepackage{float}
\usepackage{colortbl}
\usepackage{pdflscape}
\usepackage{tabularx}
\usepackage{xltabular}
\usepackage{threeparttable}
\usepackage{threeparttablex}
\usepackage[normalem]{ulem}
\usepackage{makecell}
\usepackage{xcolor}
\usepackage{float}
\usepackage{booktabs}
\usepackage{caption}
\captionsetup{font=small,labelfont=bf}
\usepackage{graphicx}
\usepackage[export]{adjustbox}
\let\oldincludegraphics\includegraphics
\renewcommand{\includegraphics}[2][]{%
  \oldincludegraphics[#1,max width=\linewidth,max height=0.9\textheight,keepaspectratio]{#2}%
}
\usepackage{fvextra}
\DefineVerbatimEnvironment{Highlighting}{Verbatim}{breaklines,fontsize=\footnotesize,commandchars=\\\{\}}
\KOMAoption{captions}{tableheading}
\makeatletter
\@ifpackageloaded{caption}{}{\usepackage{caption}}
\AtBeginDocument{%
\ifdefined\contentsname
  \renewcommand*\contentsname{Table of contents}
\else
  \newcommand\contentsname{Table of contents}
\fi
\ifdefined\listfigurename
  \renewcommand*\listfigurename{List of Figures}
\else
  \newcommand\listfigurename{List of Figures}
\fi
\ifdefined\listtablename
  \renewcommand*\listtablename{List of Tables}
\else
  \newcommand\listtablename{List of Tables}
\fi
\ifdefined\figurename
  \renewcommand*\figurename{Figure}
\else
  \newcommand\figurename{Figure}
\fi
\ifdefined\tablename
  \renewcommand*\tablename{Table}
\else
  \newcommand\tablename{Table}
\fi
}
\@ifpackageloaded{float}{}{\usepackage{float}}
\floatstyle{ruled}
\@ifundefined{c@chapter}{\newfloat{codelisting}{h}{lop}}{\newfloat{codelisting}{h}{lop}[chapter]}
\floatname{codelisting}{Listing}
\newcommand*\listoflistings{\listof{codelisting}{List of Listings}}
\makeatother
\makeatletter
\makeatother
\makeatletter
\@ifpackageloaded{caption}{}{\usepackage{caption}}
\@ifpackageloaded{subcaption}{}{\usepackage{subcaption}}
\makeatother
\usepackage{bookmark}
\IfFileExists{xurl.sty}{\usepackage{xurl}}{} % add URL line breaks if available
\urlstyle{same}
\hypersetup{
  pdftitle={Taller 3: Web Scraping, APIs y Aplicaciones Interactivas},
  pdfauthor={Grupo 10},
  colorlinks=true,
  linkcolor={blue},
  filecolor={Maroon},
  citecolor={Blue},
  urlcolor={blue},
  pdfcreator={LaTeX via pandoc}}


\title{Taller 3: Web Scraping, APIs y Aplicaciones Interactivas}
\usepackage{etoolbox}
\makeatletter
\providecommand{\subtitle}[1]{% add subtitle to \maketitle
  \apptocmd{\@title}{\par {\large #1 \par}}{}{}
}
\makeatother
\subtitle{Taller de Programación en R (2026-2), Facultad de Economía,
Universidad de los Andes}
\author{Grupo 10}
\date{2026-11-01}
\begin{document}
\maketitle

\renewcommand*\contentsname{Índice de contenido}
{
\hypersetup{linkcolor=}
\setcounter{tocdepth}{2}
\tableofcontents
}

\section{Introducción y arquitectura del Taller
3}\label{introducciuxf3n-y-arquitectura-del-taller-3}

Este informe contiene el desarrollo del Taller 3 (``Datos geoespaciales,
web scraping y herramientas de visualización'') del curso Taller de
Programación en R, en la Facultad de Economía de la Universidad de los
Andes.

El taller comprende tres ejercicios interconectados, sintetizados en la
Table~\ref{tbl-alcance-ejercicios}:

\begin{longtable}[]{@{}
  >{\raggedright\arraybackslash}p{(\linewidth - 8\tabcolsep) * \real{0.1259}}
  >{\raggedright\arraybackslash}p{(\linewidth - 8\tabcolsep) * \real{0.2222}}
  >{\centering\arraybackslash}p{(\linewidth - 8\tabcolsep) * \real{0.0667}}
  >{\centering\arraybackslash}p{(\linewidth - 8\tabcolsep) * \real{0.0963}}
  >{\raggedright\arraybackslash}p{(\linewidth - 8\tabcolsep) * \real{0.4889}}@{}}
\caption{Alcance temático y ponderación de los ejercicios del Taller
3.}\label{tbl-alcance-ejercicios}\tabularnewline
\toprule\noalign{}
\begin{minipage}[b]{\linewidth}\raggedright
Ejercicio
\end{minipage} & \begin{minipage}[b]{\linewidth}\raggedright
Módulo
\end{minipage} & \begin{minipage}[b]{\linewidth}\centering
Puntos
\end{minipage} & \begin{minipage}[b]{\linewidth}\centering
Estado
\end{minipage} & \begin{minipage}[b]{\linewidth}\raggedright
Función en el taller
\end{minipage} \\
\midrule\noalign{}
\endfirsthead
\toprule\noalign{}
\begin{minipage}[b]{\linewidth}\raggedright
Ejercicio
\end{minipage} & \begin{minipage}[b]{\linewidth}\raggedright
Módulo
\end{minipage} & \begin{minipage}[b]{\linewidth}\centering
Puntos
\end{minipage} & \begin{minipage}[b]{\linewidth}\centering
Estado
\end{minipage} & \begin{minipage}[b]{\linewidth}\raggedright
Función en el taller
\end{minipage} \\
\midrule\noalign{}
\endhead
\bottomrule\noalign{}
\endlastfoot
Ejercicio 1 & Datos geoespaciales (IDF y DANE) & 1.5 & Pendiente &
Estimación del IDF ponderado departamental y generación de mapas con
shapefiles, insumo para la pestaña 1 del aplicativo interactivo. \\
Ejercicio 2 & Web scraping y API & 2.0 & Completado & Extracción de
series diarias de IBR (Banco de la República) y SOFR (Fed de Nueva
York), validación con boletines oficiales, análisis comparativo y
consolidación de series\_financieras\_ibr\_sofr.csv. \\
Ejercicio 3 & Aplicativo interactivo en Shiny & 1.5 & Pendiente &
Tablero con tres componentes: mapa departamental del IDF (Ejercicio 1),
cálculo de devaluación implícita con las tasas obtenidas (Ejercicio 2) y
una visualización en Plotly. \\
\end{longtable}

\subsection{Flujo de trabajo del
proyecto}\label{flujo-de-trabajo-del-proyecto}

El diagrama resume el recorrido de los datos desde las fuentes externas
hasta los archivos procesados y su articulación con los demás puntos del
taller:

\begin{figure}[H]

\centering{

\includegraphics[width=14.98in,height=14.93in]{Taller_3_Grupo_10_files/figure-latex/mermaid-figure-1.png}

}

\caption{\label{fig-diagrama-flujo}Figura A. Diagrama de flujo de
extracción de datos, armonización y articulación del Taller 3.}

\end{figure}%

\begin{center}\rule{0.5\linewidth}{0.5pt}\end{center}

\section{Ejercicio 2: Integración de series macroeconómicas y
financieras}\label{ejercicio-2-integraciuxf3n-de-series-macroeconuxf3micas-y-financieras}

El objetivo de este módulo es consultar, ordenar y comparar las tasas de
interés interbancarias a un día (overnight) de Colombia y Estados Unidos
entre el 1 de enero de 2020 y el 31 de octubre de 2025.

\subsection{Metodología de extracción por entidad
emisora}\label{metodologuxeda-de-extracciuxf3n-por-entidad-emisora}

Para dar estricto cumplimiento a los requerimientos técnicos del taller,
se implementó una técnica de captura diferenciada por entidad oficial:

\begin{itemize}
\tightlist
\item
  \textbf{SOFR Overnight (Estados Unidos):} consulta programática
  automatizada a la API RESTful (formato JSON) del Federal Reserve Bank
  of New York, recuperando la serie histórica diaria y respaldando la
  respuesta cruda en disco.
\item
  \textbf{IBR Overnight (Colombia):} descarga y procesamiento directo
  sobre el árbol DOM del feed SDMX-XML del Banco de la República
  mediante técnicas de web scraping y expresiones XPath
  (\texttt{IRIBRM00} y \texttt{ER}).
\end{itemize}

\subsection{Implementación de los flujos de
extracción}\label{sec-flujos-extraccion}

Para mantener el código auditable y reproducible, cada etapa se programó
en un script independiente con operaciones secuenciales, sin
encapsulaciones innecesarias. Las siguientes secciones presentan la
arquitectura técnica y la lógica aplicada en cada caso:

\subsubsection{Flujo 1: SOFR (API REST)}

\begin{itemize}
\tightlist
\item
  \textbf{Script ejecutable:}
  \href{https://github.com/MPalma21/Talleres_de_R/blob/main/Entregable_Talleres/Taller_3/Scripts/Punto_2/01_extraer_sofr_api.R}{\texttt{01\_extraer\_sofr\_api.R}}
  \emph{(Carpeta Scripts/Punto\_2/)}
\item
  \textbf{Fuente oficial:}
  \href{https://www.newyorkfed.org/markets/reference-rates/sofr}{Federal
  Reserve Bank of New York} (Markets Data API)
\item
  \textbf{Endpoint público:}
  \url{https://markets.newyorkfed.org/api/rates/secured/sofr/search.json}
\item
  \textbf{Parámetros de consulta:} \texttt{startDate\ =\ 2020-01-01} y
  \texttt{endDate\ =\ 2025-10-31} (HTTP GET / JSON)
\item
  \textbf{Destino en raw/:} \texttt{sofr\_raw.json} (crudo) y
  \texttt{sofr\_diaria\_2020\_2025.csv} (tratado)
\end{itemize}

\begin{Shaded}
\begin{Highlighting}[]
\FunctionTok{library}\NormalTok{(tidyverse)}
\FunctionTok{library}\NormalTok{(httr)}
\FunctionTok{library}\NormalTok{(jsonlite)}
\FunctionTok{library}\NormalTok{(lubridate)}
\FunctionTok{library}\NormalTok{(here)}

\NormalTok{fecha\_inicio }\OtherTok{\textless{}{-}} \StringTok{"2020{-}01{-}01"}
\NormalTok{fecha\_fin    }\OtherTok{\textless{}{-}} \StringTok{"2025{-}10{-}31"}

\NormalTok{endpoint\_sofr }\OtherTok{\textless{}{-}} \StringTok{"https://markets.newyorkfed.org/api/rates/secured/sofr/search.json"}

\CommentTok{\# 1. Petición HTTP GET al endpoint oficial de la Fed de Nueva York}
\NormalTok{respuesta\_api }\OtherTok{\textless{}{-}} \FunctionTok{GET}\NormalTok{(}
  \AttributeTok{url =}\NormalTok{ endpoint\_sofr,}
  \AttributeTok{query =} \FunctionTok{list}\NormalTok{(}\AttributeTok{startDate =}\NormalTok{ fecha\_inicio, }\AttributeTok{endDate =}\NormalTok{ fecha\_fin),}
  \FunctionTok{user\_agent}\NormalTok{(}\StringTok{"Mozilla/5.0 (Windows NT 10.0; Win64; x64)"}\NormalTok{)}
\NormalTok{)}

\CommentTok{\# 2. Respaldo del archivo crudo en JSON}
\NormalTok{contenido\_json }\OtherTok{\textless{}{-}} \FunctionTok{content}\NormalTok{(respuesta\_api, }\AttributeTok{as =} \StringTok{"text"}\NormalTok{, }\AttributeTok{encoding =} \StringTok{"UTF{-}8"}\NormalTok{)}
\FunctionTok{writeLines}\NormalTok{(contenido\_json, }\AttributeTok{con =} \StringTok{"Outputs/Data/raw/sofr\_raw.json"}\NormalTok{)}

\CommentTok{\# 3. Tratamiento, tipado y estructuración de la serie individual}
\NormalTok{datos\_crudos }\OtherTok{\textless{}{-}} \FunctionTok{fromJSON}\NormalTok{(contenido\_json, }\AttributeTok{flatten =} \ConstantTok{TRUE}\NormalTok{)}

\NormalTok{df\_sofr }\OtherTok{\textless{}{-}}\NormalTok{ datos\_crudos}\SpecialCharTok{$}\NormalTok{refRates }\SpecialCharTok{|\textgreater{}}
  \FunctionTok{as\_tibble}\NormalTok{() }\SpecialCharTok{|\textgreater{}}
  \FunctionTok{transmute}\NormalTok{(}
    \AttributeTok{fecha     =} \FunctionTok{ymd}\NormalTok{(effectiveDate),}
    \AttributeTok{tasa\_sofr =} \FunctionTok{as.numeric}\NormalTok{(percentRate)}
\NormalTok{  ) }\SpecialCharTok{|\textgreater{}}
  \FunctionTok{filter}\NormalTok{(}
    \SpecialCharTok{!}\FunctionTok{is.na}\NormalTok{(tasa\_sofr),}
\NormalTok{    fecha }\SpecialCharTok{\textgreater{}=} \FunctionTok{ymd}\NormalTok{(fecha\_inicio),}
\NormalTok{    fecha }\SpecialCharTok{\textless{}=} \FunctionTok{ymd}\NormalTok{(fecha\_fin)}
\NormalTok{  ) }\SpecialCharTok{|\textgreater{}}
  \FunctionTok{arrange}\NormalTok{(fecha)}

\CommentTok{\# 4. Exportación de la serie tratada en CSV dentro de raw/}
\FunctionTok{write\_csv}\NormalTok{(df\_sofr, }\StringTok{"Outputs/Data/raw/sofr\_diaria\_2020\_2025.csv"}\NormalTok{)}
\end{Highlighting}
\end{Shaded}

\subsubsection{Flujo 2: IBR (Web Scraping SDMX)}

\begin{itemize}
\tightlist
\item
  \textbf{Script ejecutable:}
  \href{https://github.com/MPalma21/Talleres_de_R/blob/main/Entregable_Talleres/Taller_3/Scripts/Punto_2/02_extraer_ibr_banrep.R}{\texttt{02\_extraer\_ibr\_banrep.R}}
  \emph{(Carpeta Scripts/Punto\_2/)}
\item
  \textbf{Fuente oficial:}
  \href{https://www.banrep.gov.co/es/estadisticas/indicador-bancario-referencia-ibr}{Banco
  de la República de Colombia} (portal SUAMECA / Totoro SDMX)
\item
  \textbf{URL del feed SDMX:}
  \href{https://totoro.banrep.gov.co/nsi-jax-ws/rest/data/ESTAT,DF_IBR_DAILY_HIST,1.0/all/ALL/?startPeriod=2020&dimensionAtObservation=TIME_PERIOD&detail=full}{Feed
  SDMX Totoro SUAMECA}
\item
  \textbf{Técnica aplicada:} Web Scraping XML y extracción XPath
  (\texttt{IRIBRM00} \& \texttt{ER})
\item
  \textbf{Destino en raw/:} \texttt{ibr\_raw.xml} (crudo) y
  \texttt{ibr\_diaria\_2020\_2025.csv} (tratado)
\end{itemize}

\begin{Shaded}
\begin{Highlighting}[]
\FunctionTok{library}\NormalTok{(tidyverse)}
\FunctionTok{library}\NormalTok{(httr)}
\FunctionTok{library}\NormalTok{(xml2)}
\FunctionTok{library}\NormalTok{(lubridate)}
\FunctionTok{library}\NormalTok{(here)}

\NormalTok{fecha\_inicio }\OtherTok{\textless{}{-}} \StringTok{"2020{-}01{-}01"}
\NormalTok{fecha\_fin    }\OtherTok{\textless{}{-}} \StringTok{"2025{-}10{-}31"}

\NormalTok{url\_sdmx\_banrep }\OtherTok{\textless{}{-}} \FunctionTok{paste0}\NormalTok{(}
  \StringTok{"https://totoro.banrep.gov.co/nsi{-}jax{-}ws/rest/data/ESTAT,DF\_IBR\_DAILY\_HIST,1.0/all/ALL/"}\NormalTok{,}
  \StringTok{"?startPeriod=2020\&dimensionAtObservation=TIME\_PERIOD\&detail=full"}
\NormalTok{)}

\CommentTok{\# 1. Descarga del feed web oficial SDMX{-}XML si no se encuentra en disco}
\ControlFlowTok{if}\NormalTok{ (}\SpecialCharTok{!}\FunctionTok{file.exists}\NormalTok{(}\StringTok{"Outputs/Data/raw/ibr\_raw.xml"}\NormalTok{)) \{}
\NormalTok{  respuesta\_banrep }\OtherTok{\textless{}{-}} \FunctionTok{GET}\NormalTok{(}
    \AttributeTok{url =}\NormalTok{ url\_sdmx\_banrep,}
    \AttributeTok{config =} \FunctionTok{config}\NormalTok{(}\AttributeTok{ssl\_verifypeer =} \ConstantTok{FALSE}\NormalTok{),}
    \FunctionTok{add\_headers}\NormalTok{(}\AttributeTok{Accept =} \StringTok{"application/vnd.sdmx.genericdata+xml;version=2.1"}\NormalTok{),}
    \FunctionTok{user\_agent}\NormalTok{(}\StringTok{"Mozilla/5.0"}\NormalTok{)}
\NormalTok{  )}
  \FunctionTok{writeBin}\NormalTok{(}\FunctionTok{content}\NormalTok{(respuesta\_banrep, }\AttributeTok{as =} \StringTok{"raw"}\NormalTok{), }\StringTok{"Outputs/Data/raw/ibr\_raw.xml"}\NormalTok{)}
\NormalTok{\}}

\CommentTok{\# 2. Extracción directa con XPath sobre el DOM XML}
\NormalTok{doc\_xml }\OtherTok{\textless{}{-}} \FunctionTok{read\_xml}\NormalTok{(}\StringTok{"Outputs/Data/raw/ibr\_raw.xml"}\NormalTok{)}
\NormalTok{nodo\_serie\_on }\OtherTok{\textless{}{-}} \FunctionTok{xml\_find\_first}\NormalTok{(}
\NormalTok{  doc\_xml,}
  \StringTok{".//*[local{-}name()=\textquotesingle{}Series\textquotesingle{}][.//*[local{-}name()=\textquotesingle{}Value\textquotesingle{}][@value=\textquotesingle{}IRIBRM00\textquotesingle{}] and .//*[local{-}name()=\textquotesingle{}Value\textquotesingle{}][@value=\textquotesingle{}ER\textquotesingle{}]]"}
\NormalTok{)}

\NormalTok{nodos\_obs    }\OtherTok{\textless{}{-}} \FunctionTok{xml\_find\_all}\NormalTok{(nodo\_serie\_on, }\StringTok{".//*[local{-}name()=\textquotesingle{}Obs\textquotesingle{}]"}\NormalTok{)}
\NormalTok{fechas\_crudo }\OtherTok{\textless{}{-}} \FunctionTok{xml\_attr}\NormalTok{(}\FunctionTok{xml\_find\_all}\NormalTok{(nodos\_obs, }\StringTok{".//*[local{-}name()=\textquotesingle{}ObsDimension\textquotesingle{}]"}\NormalTok{), }\StringTok{"value"}\NormalTok{)}
\NormalTok{tasas\_crudo  }\OtherTok{\textless{}{-}} \FunctionTok{xml\_attr}\NormalTok{(}\FunctionTok{xml\_find\_all}\NormalTok{(nodos\_obs, }\StringTok{".//*[local{-}name()=\textquotesingle{}ObsValue\textquotesingle{}]"}\NormalTok{), }\StringTok{"value"}\NormalTok{)}

\CommentTok{\# 3. Tratamiento, tipado y estructuración de la serie individual}
\NormalTok{df\_ibr }\OtherTok{\textless{}{-}} \FunctionTok{tibble}\NormalTok{(}
  \AttributeTok{fecha    =} \FunctionTok{ymd}\NormalTok{(fechas\_crudo),}
  \AttributeTok{tasa\_ibr =} \FunctionTok{as.numeric}\NormalTok{(tasas\_crudo)}
\NormalTok{) }\SpecialCharTok{|\textgreater{}}
  \FunctionTok{filter}\NormalTok{(}
    \SpecialCharTok{!}\FunctionTok{is.na}\NormalTok{(tasa\_ibr),}
\NormalTok{    fecha }\SpecialCharTok{\textgreater{}=} \FunctionTok{ymd}\NormalTok{(fecha\_inicio),}
\NormalTok{    fecha }\SpecialCharTok{\textless{}=} \FunctionTok{ymd}\NormalTok{(fecha\_fin)}
\NormalTok{  ) }\SpecialCharTok{|\textgreater{}}
  \FunctionTok{arrange}\NormalTok{(fecha)}

\CommentTok{\# 4. Exportación de la serie tratada en CSV dentro de raw/}
\FunctionTok{write\_csv}\NormalTok{(df\_ibr, }\StringTok{"Outputs/Data/raw/ibr\_diaria\_2020\_2025.csv"}\NormalTok{)}
\end{Highlighting}
\end{Shaded}

\subsubsection{Flujo 3: Armonización e Integración}

\begin{itemize}
\tightlist
\item
  \textbf{Script ejecutable:}
  \href{https://github.com/MPalma21/Talleres_de_R/blob/main/Entregable_Talleres/Taller_3/Scripts/Punto_2/03_analisis_financiero_punto_2.R}{\texttt{03\_analisis\_financiero\_punto\_2.R}}
  \emph{(Carpeta Scripts/Punto\_2/)}
\item
  \textbf{Insumos de entrada:} Series tratadas en CSV desde
  \texttt{Outputs/Data/raw/}
\item
  \textbf{Procesamiento:} Lectura de CSVs, empalme continuo
  (\texttt{full\_join}), forward-fill y spread
\item
  \textbf{Tratamiento de festivos:} \emph{Forward-fill} (\texttt{fill})
  según convención del mercado monetario
\item
  \textbf{Archivo final único:}
  \texttt{Outputs/Data/processed/series\_financieras\_ibr\_sofr.csv}
\end{itemize}

\begin{Shaded}
\begin{Highlighting}[]
\FunctionTok{library}\NormalTok{(tidyverse)}
\FunctionTok{library}\NormalTok{(scales)}
\FunctionTok{library}\NormalTok{(lubridate)}
\FunctionTok{library}\NormalTok{(here)}

\CommentTok{\# 1. Lectura de las series tratadas directamente desde raw/}
\NormalTok{df\_sofr }\OtherTok{\textless{}{-}} \FunctionTok{read\_csv}\NormalTok{(}\StringTok{"Outputs/Data/raw/sofr\_diaria\_2020\_2025.csv"}\NormalTok{, }\AttributeTok{show\_col\_types =} \ConstantTok{FALSE}\NormalTok{) }\SpecialCharTok{|\textgreater{}}
  \FunctionTok{mutate}\NormalTok{(}\AttributeTok{fecha =} \FunctionTok{as.Date}\NormalTok{(fecha))}

\NormalTok{df\_ibr }\OtherTok{\textless{}{-}} \FunctionTok{read\_csv}\NormalTok{(}\StringTok{"Outputs/Data/raw/ibr\_diaria\_2020\_2025.csv"}\NormalTok{, }\AttributeTok{show\_col\_types =} \ConstantTok{FALSE}\NormalTok{) }\SpecialCharTok{|\textgreater{}}
  \FunctionTok{mutate}\NormalTok{(}\AttributeTok{fecha =} \FunctionTok{as.Date}\NormalTok{(fecha))}

\CommentTok{\# 2. Empalme temporal continuo, tratamiento de festivos (forward{-}fill) y cálculo del spread}
\NormalTok{df\_consolidado }\OtherTok{\textless{}{-}} \FunctionTok{full\_join}\NormalTok{(df\_sofr, df\_ibr, }\AttributeTok{by =} \StringTok{"fecha"}\NormalTok{) }\SpecialCharTok{|\textgreater{}}
  \FunctionTok{arrange}\NormalTok{(fecha) }\SpecialCharTok{|\textgreater{}}
  \FunctionTok{fill}\NormalTok{(tasa\_sofr, }\AttributeTok{.direction =} \StringTok{"down"}\NormalTok{) }\SpecialCharTok{|\textgreater{}}
  \FunctionTok{fill}\NormalTok{(tasa\_ibr, }\AttributeTok{.direction =} \StringTok{"down"}\NormalTok{) }\SpecialCharTok{|\textgreater{}}
  \FunctionTok{filter}\NormalTok{(}\SpecialCharTok{!}\FunctionTok{is.na}\NormalTok{(tasa\_sofr), }\SpecialCharTok{!}\FunctionTok{is.na}\NormalTok{(tasa\_ibr)) }\SpecialCharTok{|\textgreater{}}
  \FunctionTok{mutate}\NormalTok{(}\AttributeTok{spread\_ibr\_sofr =}\NormalTok{ tasa\_ibr }\SpecialCharTok{{-}}\NormalTok{ tasa\_sofr)}

\CommentTok{\# 3. Exportación del único archivo final en la carpeta processed/}
\FunctionTok{write\_csv}\NormalTok{(df\_consolidado, }\StringTok{"Outputs/Data/processed/series\_financieras\_ibr\_sofr.csv"}\NormalTok{)}
\end{Highlighting}
\end{Shaded}

\begin{center}\rule{0.5\linewidth}{0.5pt}\end{center}

\subsection{Auditoría, trazabilidad y control de
calidad}\label{auditoruxeda-trazabilidad-y-control-de-calidad}

Con el fin de garantizar la reproducibilidad y el rigor metodológico
exigidos en el taller, se implementó un protocolo de auditoría en tres
dimensiones:

\begin{enumerate}
\def\labelenumi{\arabic{enumi}.}
\tightlist
\item
  \textbf{Trazabilidad de fuentes y endpoints:} Registro sistemático de
  las entidades oficiales emisoras, métodos de extracción, scripts
  ejecutables asociados y especificación de parámetros.
\item
  \textbf{Matriz de contraste frente a boletines oficiales:}
  Confrontación puntual de tres hitos macroeconómicos clave (línea base
  pre-pandémica, pico contractivo y fecha de corte) contra las
  publicaciones oficiales del Federal Reserve Bank of New York y del
  Banco de la República.
\item
  \textbf{Evidencias de consistencia técnica:} Demostración empírica del
  tratamiento de festivos asimétricos (\emph{forward-fill}) y
  comprobación matemática de los cálculos financieros por un
  procedimiento analítico alternativo.
\end{enumerate}

\subsubsection{Trazabilidad técnica y metadatos de las fuentes
oficiales}\label{trazabilidad-tuxe9cnica-y-metadatos-de-las-fuentes-oficiales}

La Table~\ref{tbl-trazabilidad-fuentes} detalla los canales de captura,
enlaces web directos y scripts del repositorio empleados en el flujo de
extracción.

\begin{longtable}[]{@{}
  >{\raggedright\arraybackslash}p{(\linewidth - 10\tabcolsep) * \real{0.1215}}
  >{\raggedright\arraybackslash}p{(\linewidth - 10\tabcolsep) * \real{0.1776}}
  >{\raggedright\arraybackslash}p{(\linewidth - 10\tabcolsep) * \real{0.1963}}
  >{\raggedright\arraybackslash}p{(\linewidth - 10\tabcolsep) * \real{0.2150}}
  >{\centering\arraybackslash}p{(\linewidth - 10\tabcolsep) * \real{0.1121}}
  >{\raggedright\arraybackslash}p{(\linewidth - 10\tabcolsep) * \real{0.1776}}@{}}
\caption{Trazabilidad técnica y metadatos de las fuentes oficiales de
información.}\label{tbl-trazabilidad-fuentes}\tabularnewline
\toprule\noalign{}
\begin{minipage}[b]{\linewidth}\raggedright
Serie y Plazo
\end{minipage} & \begin{minipage}[b]{\linewidth}\raggedright
Emisor Oficial
\end{minipage} & \begin{minipage}[b]{\linewidth}\raggedright
Técnica de Captura y Script
\end{minipage} & \begin{minipage}[b]{\linewidth}\raggedright
Recurso Público / Endpoint
\end{minipage} & \begin{minipage}[b]{\linewidth}\centering
Observaciones
\end{minipage} & \begin{minipage}[b]{\linewidth}\raggedright
Convención Financiera
\end{minipage} \\
\midrule\noalign{}
\endfirsthead
\toprule\noalign{}
\begin{minipage}[b]{\linewidth}\raggedright
Serie y Plazo
\end{minipage} & \begin{minipage}[b]{\linewidth}\raggedright
Emisor Oficial
\end{minipage} & \begin{minipage}[b]{\linewidth}\raggedright
Técnica de Captura y Script
\end{minipage} & \begin{minipage}[b]{\linewidth}\raggedright
Recurso Público / Endpoint
\end{minipage} & \begin{minipage}[b]{\linewidth}\centering
Observaciones
\end{minipage} & \begin{minipage}[b]{\linewidth}\raggedright
Convención Financiera
\end{minipage} \\
\midrule\noalign{}
\endhead
\bottomrule\noalign{}
\endlastfoot
\textbf{SOFR Overnight} &
\href{https://www.newyorkfed.org/markets/reference-rates/sofr}{Federal
Reserve Bank of NY} & \hyperref[sec-flujos-extraccion]{API REST (JSON) -
Script 01} &
\href{https://markets.newyorkfed.org/api/rates/secured/sofr/search.json}{Markets
Data API Endpoint} & 1,458 obs. \emph{(2020-2025)} & \% Nominal Anual
(Actual/360) \\
\textbf{IBR Overnight} &
\href{https://www.banrep.gov.co/es/estadisticas/indicador-bancario-referencia-ibr}{Banco
de la República} & \hyperref[sec-flujos-extraccion]{Web Scraping XML -
Script 02} &
\href{https://totoro.banrep.gov.co/nsi-jax-ws/rest/data/ESTAT,DF_IBR_DAILY_HIST,1.0/all/ALL/?startPeriod=2020&dimensionAtObservation=TIME_PERIOD&detail=full}{Feed
SDMX-XML SUAMECA} & 1,426 obs. \emph{(2020-2025)} & \% Efectivo Anual
E.A. (Base 360) \\
\end{longtable}

\subsubsection{Matriz de contraste frente a boletines
oficiales}\label{matriz-de-contraste-frente-a-boletines-oficiales}

La Table~\ref{tbl-validacion-boletines} confronta verticalmente los
valores descargados por los scripts en R frente a los boletines
oficiales en las tres fechas de control exigidas por la guía.

\begin{longtable}[]{@{}
  >{\raggedright\arraybackslash}p{(\linewidth - 10\tabcolsep) * \real{0.1481}}
  >{\centering\arraybackslash}p{(\linewidth - 10\tabcolsep) * \real{0.1852}}
  >{\raggedright\arraybackslash}p{(\linewidth - 10\tabcolsep) * \real{0.1481}}
  >{\centering\arraybackslash}p{(\linewidth - 10\tabcolsep) * \real{0.1852}}
  >{\raggedright\arraybackslash}p{(\linewidth - 10\tabcolsep) * \real{0.1481}}
  >{\centering\arraybackslash}p{(\linewidth - 10\tabcolsep) * \real{0.1852}}@{}}
\caption{Matriz de validación y contraste con boletines oficiales en
hitos de control.}\label{tbl-validacion-boletines}\tabularnewline
\toprule\noalign{}
\begin{minipage}[b]{\linewidth}\raggedright
Hito Macroeconómico
\end{minipage} & \begin{minipage}[b]{\linewidth}\centering
Fecha
\end{minipage} & \begin{minipage}[b]{\linewidth}\raggedright
Serie Evaluada
\end{minipage} & \begin{minipage}[b]{\linewidth}\centering
Tasa en R
\end{minipage} & \begin{minipage}[b]{\linewidth}\raggedright
Boletín Oficial de Referencia
\end{minipage} & \begin{minipage}[b]{\linewidth}\centering
Estado de Auditoría
\end{minipage} \\
\midrule\noalign{}
\endfirsthead
\toprule\noalign{}
\begin{minipage}[b]{\linewidth}\raggedright
Hito Macroeconómico
\end{minipage} & \begin{minipage}[b]{\linewidth}\centering
Fecha
\end{minipage} & \begin{minipage}[b]{\linewidth}\raggedright
Serie Evaluada
\end{minipage} & \begin{minipage}[b]{\linewidth}\centering
Tasa en R
\end{minipage} & \begin{minipage}[b]{\linewidth}\raggedright
Boletín Oficial de Referencia
\end{minipage} & \begin{minipage}[b]{\linewidth}\centering
Estado de Auditoría
\end{minipage} \\
\midrule\noalign{}
\endhead
\bottomrule\noalign{}
\endlastfoot
\textbf{1. Línea base pre-pandémica} & \texttt{2020-01-02} & SOFR ON
(EE. UU.) & \textbf{1.54\%} &
\href{https://markets.newyorkfed.org/api/rates/secured/sofr/search.json?startDate=2020-01-02&endDate=2020-01-02}{1.54\%
(NY Fed Markets API)} & {Coincidencia Exacta} \\
\textbf{1. Línea base pre-pandémica} & \texttt{2020-01-02} & IBR ON
(Colombia) & \textbf{4.259\%} &
\href{https://totoro.banrep.gov.co/nsi-jax-ws/rest/data/ESTAT,DF_IBR_DAILY_HIST,1.0/all/ALL/?startPeriod=2020-01-02&endPeriod=2020-01-02}{4.259\%
E.A. (BanRep SUAMECA)} & {Coincidencia Exacta} \\
\textbf{2. Pico contractivo antiinflacionario} & \texttt{2023-07-27} &
SOFR ON (EE. UU.) & \textbf{5.30\%} &
\href{https://markets.newyorkfed.org/api/rates/secured/sofr/search.json?startDate=2023-07-27&endDate=2023-07-27}{5.30\%
(Corredor Fed 5.25\% - 5.50\%)} & {Coincidencia Exacta} \\
\textbf{2. Pico contractivo antiinflacionario} & \texttt{2023-04-28} &
IBR ON (Colombia) & \textbf{13.250\%} &
\href{https://totoro.banrep.gov.co/nsi-jax-ws/rest/data/ESTAT,DF_IBR_DAILY_HIST,1.0/all/ALL/?startPeriod=2023-04-28&endPeriod=2023-04-28}{13.250\%
E.A. (Repo BanRep 13.25\%)} & {Coincidencia Exacta} \\
\textbf{3. Cierre del horizonte evaluado} & \texttt{2025-10-31} & SOFR
ON (EE. UU.) & \textbf{4.22\%} &
\href{https://markets.newyorkfed.org/api/rates/secured/sofr/search.json?startDate=2025-10-31&endDate=2025-10-31}{4.22\%
(NY Fed Markets API)} & {Coincidencia Exacta} \\
\textbf{3. Cierre del horizonte evaluado} & \texttt{2025-10-31} & IBR ON
(Colombia) & \textbf{9.250\%} &
\href{https://totoro.banrep.gov.co/nsi-jax-ws/rest/data/ESTAT,DF_IBR_DAILY_HIST,1.0/all/ALL/?startPeriod=2025-10-31&endPeriod=2025-10-31}{9.250\%
E.A. (BanRep SUAMECA)} & {Coincidencia Exacta} \\
\end{longtable}

\subsubsection{Evidencias de auditoría
metodológica}\label{evidencias-de-auditoruxeda-metodoluxf3gica}

\paragraph{Evidencia 1: Tratamiento empírico de festivos
asimétricos}\label{evidencia-1-tratamiento-empuxedrico-de-festivos-asimuxe9tricos}

La disparidad de calendarios bancarios entre Colombia (festivos por Ley
Emiliani y celebraciones religiosas) y Estados Unidos (festivos
federales) genera desalineaciones temporales. La
Table~\ref{tbl-evidencia-festivos} expone casos concretos donde la
imputación mediante \emph{forward-fill}
(\texttt{fill(...,\ .direction\ =\ "down")}) garantizó la continuidad
del devengo financiero sin alterar la dinámica de tasas.

\begin{longtable}[]{@{}
  >{\centering\arraybackslash}p{(\linewidth - 8\tabcolsep) * \real{0.2381}}
  >{\raggedright\arraybackslash}p{(\linewidth - 8\tabcolsep) * \real{0.1905}}
  >{\raggedright\arraybackslash}p{(\linewidth - 8\tabcolsep) * \real{0.1905}}
  >{\raggedright\arraybackslash}p{(\linewidth - 8\tabcolsep) * \real{0.1905}}
  >{\raggedright\arraybackslash}p{(\linewidth - 8\tabcolsep) * \real{0.1905}}@{}}
\caption{Casos testigo de armonización de festivos no coincidentes entre
Colombia y EE. UU.}\label{tbl-evidencia-festivos}\tabularnewline
\toprule\noalign{}
\begin{minipage}[b]{\linewidth}\centering
Fecha
\end{minipage} & \begin{minipage}[b]{\linewidth}\raggedright
Contexto de Mercado
\end{minipage} & \begin{minipage}[b]{\linewidth}\raggedright
Dato Crudo Extraído
\end{minipage} & \begin{minipage}[b]{\linewidth}\raggedright
Tratamiento Aplicado
\end{minipage} & \begin{minipage}[b]{\linewidth}\raggedright
Justificación Financiera
\end{minipage} \\
\midrule\noalign{}
\endfirsthead
\toprule\noalign{}
\begin{minipage}[b]{\linewidth}\centering
Fecha
\end{minipage} & \begin{minipage}[b]{\linewidth}\raggedright
Contexto de Mercado
\end{minipage} & \begin{minipage}[b]{\linewidth}\raggedright
Dato Crudo Extraído
\end{minipage} & \begin{minipage}[b]{\linewidth}\raggedright
Tratamiento Aplicado
\end{minipage} & \begin{minipage}[b]{\linewidth}\raggedright
Justificación Financiera
\end{minipage} \\
\midrule\noalign{}
\endhead
\bottomrule\noalign{}
\endlastfoot
\texttt{2020-01-06} & \textbf{Día de Reyes Magos} (Festivo Colombia,
hábil EE. UU.) & SOFR: 1.55\%, IBR: \texttt{NA} & Forward-fill: IBR
mantiene \textbf{4.259\%} (del 2020-01-03) & Los contratos monetarios
overnight continúan devengando intereses a la última tasa oficial
vigente hasta la reapertura bancaria. \\
\texttt{2020-01-20} & \textbf{Martin Luther King Jr.~Day} (Festivo EE.
UU., hábil Colombia) & SOFR: \texttt{NA}, IBR: 4.250\% & Forward-fill:
SOFR mantiene \textbf{1.55\%} (del 2020-01-17) & Convención bancaria
internacional estándar (\emph{last observation carried forward}). \\
\texttt{2024-07-04} & \textbf{Independence Day} (Festivo EE. UU., hábil
Colombia) & SOFR: \texttt{NA}, IBR: 11.230\% & Forward-fill: SOFR
mantiene \textbf{5.34\%} (del 2024-07-03) & Conserva la integridad
temporal de la base continua (1,500 observaciones de mercado). \\
\texttt{2024-07-20} & \textbf{Grito de Independencia} (Festivo Colombia,
hábil EE. UU.) & SOFR: 5.34\%, IBR: \texttt{NA} & Forward-fill: IBR
mantiene \textbf{11.230\%} (del 2024-07-19) & Evita truncamientos
artificiales en el vector diario del spread financiero. \\
\end{longtable}

\paragraph{Evidencia 2: Comprobación analítica por procedimiento
alternativo}\label{evidencia-2-comprobaciuxf3n-analuxedtica-por-procedimiento-alternativo}

Para contrastar la exactitud matemática de las transformaciones
ejecutadas en R, se verificaron analíticamente el spread financiero y la
devaluación implícita para la fecha de corte (\textbf{31 de octubre de
2025}):

\begin{enumerate}
\def\labelenumi{\arabic{enumi}.}
\item
  \textbf{Parámetros de entrada:}

  \begin{itemize}
  \tightlist
  \item
    Tasa doméstica: \(i_d = \text{IBR}_{ON} = 9.250\% = 0.09250\).
  \item
    Tasa externa: \(i_e = \text{SOFR}_{ON} = 4.220\% = 0.04220\).
  \end{itemize}
\item
  \textbf{Cálculo del spread financiero:}
  \[\text{Spread}_t = i_d - i_e = 9.250\% - 4.220\% = 5.030\% \quad (503.0 \text{ puntos básicos})\]
\item
  \textbf{Cálculo exacto de la devaluación implícita (\(\pi\)):}
  \[\begin{aligned}
  \pi &= \frac{1 + i_d}{1 + i_e} - 1 = \frac{1 + 0.09250}{1 + 0.04220} - 1 \\
  &= \frac{1.09250}{1.04220} - 1 \approx 0.048263 \implies 4.826\%
  \end{aligned}\]
\item
  \textbf{Procedimiento alternativo (Expansión de Taylor de primer
  orden):} \[\begin{aligned}
  \pi &\approx (i_d - i_e) - i_e(i_d - i_e) \\
  &= 0.05030 - (0.04220 \times 0.05030) \\
  &= 0.05030 - 0.002123 = 0.048177 \implies 4.818\%
  \end{aligned}\] La diferencia residual de 0.008 puntos porcentuales (8
  puntos básicos) corresponde estrictamente al término no lineal de
  segundo orden (\(i_e \cdot \pi\)), validando la consistencia
  computacional del flujo programado.
\end{enumerate}

\begin{center}\rule{0.5\linewidth}{0.5pt}\end{center}

\subsection{Literal 2.A: Gráfico y análisis del IBR Overnight
(2020-2025)}\label{literal-2.a-gruxe1fico-y-anuxe1lisis-del-ibr-overnight-2020-2025}

\subsubsection{Unidad de medida y convención
financiera}\label{unidad-de-medida-y-convenciuxf3n-financiera}

El Indicador Bancario de Referencia (IBR) a plazo overnight mide el
costo diario al que los bancos comerciales colombianos están dispuestos
a prestar o captar liquidez en pesos.

\begin{itemize}
\tightlist
\item
  Unidad: Tasa Efectiva Anual (E.A., en porcentaje).
\item
  Convención de conteo: Base comercial de 360 días (doce meses de 30
  días, estándar del sistema financiero colombiano).
\item
  Relación entre tasa nominal y efectiva:
  \[i_{E.A.} = \left( 1 + \frac{i_{nominal} \times 1}{360} \right)^{360} - 1\]
\end{itemize}

\subsubsection{Visualización
histórica}\label{visualizaciuxf3n-histuxf3rica}

\begin{figure}[H]

\centering{

\pandocbounded{\includegraphics[keepaspectratio]{Taller_3_Grupo_10_files/figure-pdf/fig-ibr-diaria-1.pdf}}

}

\caption{\label{fig-ibr-diaria}Figura 1. Evolución diaria del IBR
Overnight en Colombia (2020-2025), señalando el estímulo inicial por
COVID-19, el ajuste alcista y el período de normalización.}

\end{figure}%

\subsubsection{Comportamiento de la
serie}\label{comportamiento-de-la-serie}

\paragraph{Estímulo monetario por COVID-19 (marzo de 2020 a septiembre
de
2021)}\label{estuxedmulo-monetario-por-covid-19-marzo-de-2020-a-septiembre-de-2021}

La llegada de la pandemia en marzo de 2020 llevó a la Junta Directiva
del Banco de la República a recortar su tasa de intervención en siete
reuniones consecutivas, pasando de 4.25\% a un mínimo histórico de
1.75\% a finales de septiembre de 2020. El IBR overnight transmitió este
impulso con rapidez y se mantuvo en una banda estrecha entre 1.70\% y
1.75\% durante doce meses, abaratando el costo del fondeo a muy corto
plazo para los intermediarios financieros en una etapa de parálisis
productiva.

Este piso de tasas estuvo respaldado por operaciones transitorias de
liquidez (repos) en títulos públicos y privados, además de compras
definitivas de títulos para evitar tensiones de liquidez en el sistema
de pagos. La permanencia del IBR en 1.75\% evidenció un exceso sostenido
de reservas bancarias, donde los establecimientos de crédito mantuvieron
saldos disponibles suficientes sin necesidad de pujar al alza por
recursos interbancarios.

\paragraph{Choque inflacionario y ciclo alcista (octubre de 2021 a abril
de
2023)}\label{choque-inflacionario-y-ciclo-alcista-octubre-de-2021-a-abril-de-2023}

Hacia el último trimestre de 2021, la reapertura comercial, los cuellos
de botella en el comercio internacional y el encarecimiento de los
insumos reactivaron las presiones sobre los precios al consumidor. El
fenómeno se profundizó durante 2022 por el choque de materias primas
derivado del conflicto en Ucrania y la depreciación del peso, lo que
llevó la inflación anual por encima del 13\% en los primeros meses de
2023. Ante el riesgo de desanclar las expectativas de inflación de
mediano plazo, el emisor inició el ciclo contractivo más acelerado en la
historia del esquema de inflación objetivo.

Entre octubre de 2021 y abril de 2023, la tasa de intervención subió de
forma acumulada 1,150 puntos básicos, desde 1.75\% hasta 13.25\%. La
serie diaria del IBR refleja este proceso como una escalera ascendente
continua que replica de forma casi inmediata cada incremento de la
Junta, alcanzando su nivel máximo el 28 de abril de 2023. El ajuste
incrementó el costo del crédito en todas las modalidades y restringió la
demanda interna para encauzar la inflación hacia la meta.

\paragraph{Meseta restrictiva y normalización paulatina (mayo de 2023 a
octubre de
2025)}\label{meseta-restrictiva-y-normalizaciuxf3n-paulatina-mayo-de-2023-a-octubre-de-2025}

Alcanzado el techo de 13.25\% en abril de 2023, la Junta mantuvo la tasa
sin cambios durante siete meses consecutivos. Esta pausa prolongada
buscaba verificar que la desaceleración del índice de precios fuera
consistente antes de iniciar cualquier relajamiento, dada la rigidez que
mostraron los rubros de servicios y arriendos por los mecanismos
formales de indexación.

En diciembre de 2023, con la inflación anual descendiendo y la inversión
mostrando caídas continuas, el banco central comenzó a recortar la tasa
directriz en pasos sucesivos de 50 puntos básicos. Esta fase de
flexibilización situó el IBR overnight en 9.25\% al cierre de octubre de
2025. Aunque el nivel actual es sustancialmente menor que el pico de
2023, la tasa real permanece en terreno neutral o moderadamente
restrictivo, lo que busca afianzar la convergencia gradual de los
precios hacia la meta del 3\%.

\begin{center}\rule{0.5\linewidth}{0.5pt}\end{center}

\subsection{Literal 2.B: Gráfico y análisis del SOFR Overnight
(2020-2025)}\label{literal-2.b-gruxe1fico-y-anuxe1lisis-del-sofr-overnight-2020-2025}

\subsubsection{Unidad de medida y convención
financiera}\label{unidad-de-medida-y-convenciuxf3n-financiera-1}

La tasa SOFR (Secured Overnight Financing Rate) mide el costo del fondeo
en dólares a 24 horas garantizado con títulos del Tesoro de Estados
Unidos en el mercado repo tripartito.

\begin{itemize}
\tightlist
\item
  Unidad: Tasa porcentual anualizada nominal (APR, \%).
\item
  Convención de conteo: Base del mercado monetario estadounidense
  Actual/360 (días calendario reales sobre 360).
\item
  Respaldo colateral: al negociarse contra deuda soberana de Estados
  Unidos, la tasa no incorpora primas de riesgo crediticio de
  contraparte bancaria, a diferencia del extinto esquema LIBOR.
\end{itemize}

\subsubsection{Visualización
histórica}\label{visualizaciuxf3n-histuxf3rica-1}

\begin{figure}[H]

\centering{

\pandocbounded{\includegraphics[keepaspectratio]{Taller_3_Grupo_10_files/figure-pdf/fig-sofr-diaria-1.pdf}}

}

\caption{\label{fig-sofr-diaria}Figura 2. Evolución diaria del SOFR
Overnight en Estados Unidos (2020-2025), mostrando el período de tasa
cero, el ciclo de aumentos de la Reserva Federal y los recortes
posteriores.}

\end{figure}%

\subsubsection{Comportamiento de la
serie}\label{comportamiento-de-la-serie-1}

\paragraph{Política de tasa cero (ZIRP, marzo de 2020 a marzo de
2022)}\label{poluxedtica-de-tasa-cero-zirp-marzo-de-2020-a-marzo-de-2022}

Para contener el colapso financiero al inicio de la pandemia en marzo de
2020, el Comité Federal de Mercado Abierto (FOMC) redujo su tasa de
referencia al rango de 0.00\% a 0.25\% e inició compras a gran escala de
bonos del Tesoro y títulos respaldados por hipotecas. La tasa SOFR cayó
rápidamente desde 1.54\% a principios de 2020 hasta situarse muy cerca
de cero, promediando 0.05\% a lo largo de dos años continuos.

La gran cantidad de reservas inyectadas al sistema bancario generó
presiones a la baja sobre las tasas a un día. Para evitar que el costo
del fondeo colateralizado pasara a terreno negativo, la Reserva Federal
utilizó la facilidad de operaciones de recompra inversa (Overnight
Reverse Repo Facility, ON RRP), que funcionó como un piso operativo en
torno a 0.05\%. Este mecanismo permitió absorber el excedente de
liquidez y mantener la estabilidad del mercado repo sin las tensiones
observadas en septiembre de 2019.

\paragraph{Choque inflacionario y endurecimiento acelerado (marzo de
2022 a julio de
2023)}\label{choque-inflacionario-y-endurecimiento-acelerado-marzo-de-2022-a-julio-de-2023}

En los primeros meses de 2022, la combinación de problemas en las
cadenas globales de valor, un mercado laboral con desempleo en mínimos
históricos y el estímulo fiscal pospandemia elevó la inflación anual al
9.1\% en junio de 2022, el nivel más alto en cuarenta años. La Reserva
Federal abandonó la premisa de transitoriedad e implementó una de las
correcciones monetarias más rápidas registradas en las últimas décadas.

A partir de marzo de 2022, el FOMC aplicó incrementos sucesivos que
incluyeron cuatro movimientos consecutivos de 75 puntos básicos,
situando el corredor de política en 5.25\% - 5.50\% en julio de 2023. La
serie diaria del SOFR siguió de cerca cada decisión del banco central
estadounidense y escaló desde 0.08\% hasta marcar un pico de 5.30\% el
27 de julio de 2023. La evolución de la tasa ilustra la eficiencia con
la que los ajustes en las facilidades de la Fed se transmiten al mercado
de deuda soberana de corto plazo.

\paragraph{Meseta prolongada y ciclo de recortes (agosto de 2023 a
octubre de
2025)}\label{meseta-prolongada-y-ciclo-de-recortes-agosto-de-2023-a-octubre-de-2025}

Entre agosto de 2023 y agosto de 2024, la Reserva Federal mantuvo su
tasa de intervención inalterada bajo la directriz de sostener
condiciones restrictivas el tiempo necesario para doblegar la inflación.
Durante este año de meseta, el SOFR osciló en una franja muy estable
entre 5.30\% y 5.33\%, con fluctuaciones menores explicadas
principalmente por efectos de calendario y fechas de liquidación del
Tesoro.

En septiembre de 2024, con la inflación subyacente acercándose a la meta
del 2\% y señales de enfriamiento en el empleo, la Fed comenzó a
flexibilizar su postura con un recorte inicial de 50 puntos básicos,
seguido de reducciones adicionales de 25 puntos básicos. Hacia el 31 de
octubre de 2025, el SOFR cerró en 4.22\%, reflejando la transición hacia
un entorno monetario menos restrictivo destinado a preservar el
crecimiento sin descuidar el equilibrio de precios.

\begin{center}\rule{0.5\linewidth}{0.5pt}\end{center}

\subsection{Literal 2.C: Gráfico con doble eje y análisis
comparativo}\label{literal-2.c-gruxe1fico-con-doble-eje-y-anuxe1lisis-comparativo}

\subsubsection{Justificación de la frecuencia
diaria}\label{justificaciuxf3n-de-la-frecuencia-diaria}

Los mercados monetarios overnight operan con vencimiento a 24 horas. El
uso de datos diarios resulta indispensable para observar la volatilidad
en cierres de mes, pagos tributarios y subastas de liquidez, efectos que
un promedio mensual o trimestral borraría por completo.

\subsubsection{Tratamiento de días no bursátiles
cruzados}\label{tratamiento-de-duxedas-no-bursuxe1tiles-cruzados}

Los calendarios festivos de ambos países no coinciden: Colombia cuenta
con los feriados de la Ley Emiliani y fechas religiosas, mientras que
Estados Unidos observa festivos federales como el Día de Martin Luther
King o el 4 de Julio.

Para integrar ambas series en una misma serie temporal sin distorsiones,
se combinaron las observaciones mediante \texttt{full\_join} y se imputó
el valor inmediato anterior
(\texttt{fill(...,\ .direction\ =\ "down")}). Esta regla refleja la
práctica financiera habitual: cuando un mercado cierra por feriado
local, los contratos vigentes continúan acumulando intereses a la última
tasa oficial fijada hasta la siguiente reapertura hábil.

\subsubsection{Visualización conjunta con doble
eje}\label{visualizaciuxf3n-conjunta-con-doble-eje}

\begin{figure}[H]

\centering{

\pandocbounded{\includegraphics[keepaspectratio]{Taller_3_Grupo_10_files/figure-pdf/fig-doble-eje-1.pdf}}

}

\caption{\label{fig-doble-eje}Figura 3. Comparación de tasas de
referencia: SOFR en el eje izquierdo e IBR en el eje derecho
(2020-2025).}

\end{figure}%

\subsubsection{Interpretación
económica}\label{interpretaciuxf3n-econuxf3mica}

\paragraph{Paridad descubierta de tasas y prima de
riesgo}\label{paridad-descubierta-de-tasas-y-prima-de-riesgo}

La diferencia en la magnitud de los ciclos (el SOFR subió 525 puntos
básicos, mientras que el IBR aumentó 1,150) ilustra la dinámica de una
economía emergente frente al emisor de la moneda de reserva global. Con
movilidad de capitales y régimen cambiario flexible, la paridad
descubierta de tasas de interés (UIP) aproximada incorpora la prima por
riesgo soberano:

\[i_{t}^{COL} \approx i_{t}^{USA} + \mathbb{E}_t[\Delta s_{t+1}] + \rho_t\]

En esta expresión, \(s_t\) representa el logaritmo del tipo de cambio
peso/dólar y \(\rho_t\) resume el riesgo país (medido comúnmente por el
EMBI) y la prima por inflación. Durante 2022 y 2023, el aumento en la
aversión global al riesgo y la depreciación del peso obligaron al Banco
de la República a fijar un diferencial de tasas sustancialmente más alto
para estabilizar las presiones cambiarias y evitar salidas desordenadas
de capital.

\paragraph{Evolución del spread financiero (IBR -
SOFR)}\label{evoluciuxf3n-del-spread-financiero-ibr---sofr}

Entre 2020 y 2021, el diferencial entre ambas tasas fluctuó entre 150 y
200 puntos básicos, un margen moderado que respondió a la abundancia de
liquidez global y a la reducción sincrónica del costo del dinero por la
emergencia sanitaria.

A partir de 2022, la mayor persistencia inflacionaria en Colombia y el
encarecimiento de la deuda externa ampliaron este diferencial de forma
considerable, superando los 800 puntos básicos durante el primer
semestre de 2023 cuando el IBR alcanzó el 13.25\% frente a un SOFR de
5.05\%. Con el avance del proceso desinflacionario y el inicio de los
recortes en ambos emisores, el spread se comprimió paulatinamente hasta
ubicarse en 503 puntos básicos al 31 de octubre de 2025 (9.25\% frente a
4.22\%), reflejando una normalización relativa de las condiciones
monetarias comparadas.

\paragraph{Rezagos en la transmisión
monetaria}\label{rezagos-en-la-transmisiuxf3n-monetaria}

Como observó Milton Friedman, las decisiones de política monetaria
operan con rezagos temporales variables. El Banco de la República inició
su ciclo alcista en octubre de 2021, unos seis meses antes que la
Reserva Federal (marzo de 2022), lo que concuerda con la mayor
sensibilidad inicial de los precios en economías en desarrollo ante
choques de alimentos y materias primas. De forma similar, la
persistencia de la inflación local obligó a mantener niveles
restrictivos por un período más largo que el de las economías avanzadas.

\begin{center}\rule{0.5\linewidth}{0.5pt}\end{center}

\section{Conclusiones}\label{conclusiones}

\begin{enumerate}
\def\labelenumi{\arabic{enumi}.}
\tightlist
\item
  Los flujos de captura se ejecutaron mediante una API REST para el SOFR
  y web scraping sobre documentos SDMX-XML para el IBR, cumpliendo el
  requerimiento de emplear metodologías diferenciadas por entidad.
\item
  El desfase entre los días festivos de Colombia y Estados Unidos se
  resolvió mediante un calendario unificado con arrastre del último
  valor observado (\emph{forward-fill}), lo que garantiza la
  comparabilidad diaria del diferencial de tasas.
\item
  La base continua consolidada
  (\texttt{series\_financieras\_ibr\_sofr.csv}) queda almacenada en
  \texttt{Outputs/Data/processed/}, lista para alimentar la calculadora
  de devaluación implícita en la pestaña 2 del aplicativo interactivo
  del Ejercicio 3.
\end{enumerate}




\end{document}
